% report/preamble.tex
% Visual identity for the PoC final report.
% Engine: LuaLaTeX (via `latexmk -lualatex`, see report/latexmkrc). Needs real
% OpenType fonts (Libertinus, Source Sans Pro), so pdflatex is not supported.

% ---------------------------------------------------------------- geometry
\usepackage{geometry}
\geometry{
  a4paper,
  left=28mm, right=32mm,
  top=27mm, bottom=30mm,
  marginparwidth=16mm, marginparsep=4mm,
  headsep=8mm, footskip=12mm,
}

% ------------------------------------------------------------------ fonts
\usepackage{fontspec}
\setmainfont{Libertinus Serif}
\setsansfont{Source Sans Pro}
\setmonofont{Libertinus Mono}[Scale=0.92]
\newfontfamily\semiboldfont{Source Sans Pro Semibold}

\usepackage[american]{babel}
\usepackage{microtype}
\usepackage{csquotes}

% ------------------------------------------------------------------ color
\usepackage{xcolor}
% report/style/colors.tex
% Single source of truth for the report's color palette.
% Included by report/preamble.tex (document) and report/figures/src/figstyle.tex
% (TikZ figures), so prose and figures never drift apart.
%
% Assumes xcolor is already loaded by the including file.

% --- Neutral scale ---------------------------------------------------------
\definecolor{ink}{HTML}{1A1A1A}      % body text
\definecolor{gray700}{HTML}{4D4D4D}  % captions, secondary text, footers
\definecolor{gray400}{HTML}{A6A6A6}  % rules, borders, artefact/data-store nodes
\definecolor{gray200}{HTML}{DDDCDA}  % hairlines
\definecolor{gray100}{HTML}{F4F2F0}  % subtle box/table-row tints
\definecolor{paper}{HTML}{FFFFFF}

% --- Primary accent (headings, links, Finding/Key-question boxes) ---------
\definecolor{accent}{HTML}{7A2048}       % deep burgundy
\definecolor{accentlight}{HTML}{F2E6EC}  % 8-10%-equivalent tint for fills

% --- Evidence-category palette (color-blind-safe; Okabe-Ito derived, ------
% --- darkened for AA contrast with white badge text) ----------------------
\definecolor{evlitcol}{HTML}{0B5FA5}  % ESTABLISHED LITERATURE  - blue
\definecolor{evobscol}{HTML}{00815D}  % OBSERVED IN THE PoC     - bluish green
\definecolor{evintcol}{HTML}{B36A00}  % INTERPRETATION          - amber/ochre
\definecolor{evhypcol}{HTML}{9C3D72}  % HYPOTHESIS              - purple
\definecolor{evfutcol}{HTML}{5A5A5A}  % FUTURE WORK             - neutral gray


% --------------------------------------------------------------- KOMA base
% scrartcl gives \part/\section/.../\paragraph with no chapter level, which
% matches "numbered parts and sections" without book-length chapter scaffolding.
\KOMAoptions{
  fontsize=11pt,
  parskip=half-,
  numbers=noenddot,
  headings=big,
}

\setkomafont{disposition}{\normalfont\sffamily\bfseries}
\addtokomafont{part}{\color{accent}}
\addtokomafont{partnumber}{\color{accent}}
\addtokomafont{partentry}{\color{accent}}
\addtokomafont{section}{\color{accent}}
\addtokomafont{subsection}{\color{ink}}
\addtokomafont{subsubsection}{\color{ink}}
\addtokomafont{paragraph}{\color{ink}}
\setkomafont{captionlabel}{\normalfont\sffamily\bfseries\color{gray700}}
\setkomafont{caption}{\normalfont\sffamily\small\color{gray700}}
\renewcaptionname{american}{\figurename}{Figure}
\renewcaptionname{american}{\tablename}{Table}

% Running header/footer: minimal, sans, gray.
\usepackage[automark]{scrlayer-scrpage}
\clearpairofpagestyles
\ohead{\sffamily\small\color{gray700}\headmark}
\ofoot{\sffamily\small\color{gray700}\pagemark}
\ifoot{\sffamily\small\color{gray700}Finding What Experts Leave Unsaid --- Extended Technical Report}
\setkomafont{pageheadfoot}{\normalfont\sffamily\small\color{gray700}}
\setkomafont{pagenumber}{\normalfont\sffamily\color{gray700}}

% ------------------------------------------------------------------ tables
\usepackage{booktabs}
\usepackage{array}
\usepackage{tabularx}
\usepackage{longtable}
\usepackage{placeins}
\usepackage{multirow}
\usepackage{siunitx}
\sisetup{detect-all}
% Deliberately no global round-mode/round-precision: siunitx would then
% reformat every number to a fixed decimal count regardless of table-format,
% turning plain counts (e.g. "3") into "3.00". Set rounding per table instead.

% ------------------------------------------------------------------- lists
\usepackage{enumitem}
\setlist{topsep=4pt, itemsep=2pt, parsep=0pt}

% ------------------------------------------------------------------ figures
\usepackage{graphicx}
\usepackage{caption}
\usepackage{subcaption}
\captionsetup{
  labelsep=colon,
  font={sf,small,color=gray700},
  labelfont={sf,bf,color=gray700},
  belowskip=2pt, aboveskip=6pt,
}
\usepackage{marginnote}

% ------------------------------------------------------------------ boxes
\usepackage[most]{tcolorbox}
\tcbuselibrary{skins,breakable}

% ---------------------------------------------------------------- bibliography
\usepackage[
  backend=biber,
  style=authoryear-comp,
  sorting=nyt,
  maxcitenames=2,
  uniquename=false,
  uniquelist=false,
  giveninits=true,
  doi=true,
  url=true,
  isbn=false,
  eprint=false,
]{biblatex}

% -------------------------------------------------------------------- links
\usepackage{hyperref}
\hypersetup{
  colorlinks=true,
  linkcolor=accent!80!black,
  citecolor=evlitcol,
  urlcolor=gray700,
  filecolor=gray700,
  linktoc=all,
  bookmarksnumbered=true,
  pdfstartview={XYZ null null 1},
}
\usepackage{cleveref}
\crefname{figure}{Figure}{Figures}
\crefname{table}{Table}{Tables}
\crefname{section}{Section}{Sections}
\crefname{part}{Part}{Parts}

% ============================================================================
%  EVIDENCE VISUAL LANGUAGE
%  Five statement categories, plus Finding and Key-question boxes, plus
%  inline badge/tag macros. See report/notes/10_visual_language.md for the
%  full rulebook.
% ============================================================================

% --- inline badge + small-caps word: color + letter + word (three-channel
%     redundancy so the category survives grayscale print and color blindness)
\newcommand{\evbadge}[3]{%
  \tcbox[on line, colback=#2, colframe=#2, boxrule=0pt, arc=1pt,
    left=2pt, right=2pt, top=0.3pt, bottom=0.3pt,
    fontupper=\sffamily\bfseries\fontsize{6.3}{6.3}\selectfont\color{white}]{#1}%
  \,\textsc{\sffamily\footnotesize\color{#2}#3}%
}
\newcommand{\evlit}{\evbadge{L}{evlitcol}{lit.}}
\newcommand{\evobs}{\evbadge{O}{evobscol}{obs.}}
\newcommand{\evint}{\evbadge{I}{evintcol}{int.}}
\newcommand{\evhyp}{\evbadge{H}{evhypcol}{hyp.}}
\newcommand{\evfut}{\evbadge{F}{evfutcol}{fut.}}

% --- block boxes: thin left rule only, no gradient, no full frame ----------
\tcbset{
  evbox/.style={
    enhanced, breakable,
    colback=white, coltext=ink,
    boxrule=0pt, leftrule=2.6pt,
    toprule=0pt, bottomrule=0pt, rightrule=0pt,
    arc=0pt, outer arc=0pt,
    left=9pt, right=9pt, top=6pt, bottom=6pt,
    before skip=8pt, after skip=8pt,
    fonttitle=\sffamily\bfseries\small,
    attach title to upper={\par\vspace{2pt}\noindent},
  }
}

% NB: coltitle defaults to a light color meant for a solid title bar; since
% "attach title to upper" puts the title directly on the (near-white) box
% body, each box must set coltitle explicitly or the label text is nearly
% invisible.
\newtcolorbox{evlitbox}[1][]{evbox, colframe=evlitcol, colback=evlitcol!5!white,
  coltitle=evlitcol, title={\evlit\ \textsc{Established literature}}, #1}
\newtcolorbox{evobsbox}[1][]{evbox, colframe=evobscol, colback=evobscol!5!white,
  coltitle=evobscol, title={\evobs\ \textsc{Observed in the PoC}}, #1}
\newtcolorbox{evintbox}[1][]{evbox, colframe=evintcol, colback=evintcol!5!white,
  coltitle=evintcol, title={\evint\ \textsc{Interpretation}}, #1}
\newtcolorbox{evhypbox}[1][]{evbox, colframe=evhypcol, colback=evhypcol!5!white,
  coltitle=evhypcol, title={\evhyp\ \textsc{Hypothesis}}, #1}
\newtcolorbox{evfutbox}[1][]{evbox, colframe=evfutcol, colback=evfutcol!5!white,
  leftrule=2.6pt, borderline west={2.6pt}{0pt}{evfutcol, dashed},
  coltitle=evfutcol, title={\evfut\ \textsc{Future work}}, #1}

% --- Finding: same restrained left-rule style, primary accent color -------
\newtcolorbox{findingbox}[1][]{evbox, colframe=accent, colback=accentlight,
  leftrule=3.2pt,
  title={\tcbox[on line, colback=accent, colframe=accent, boxrule=0pt, arc=1pt,
    left=2pt, right=2pt, top=0.3pt, bottom=0.3pt,
    fontupper=\sffamily\bfseries\fontsize{6.3}{6.3}\selectfont\color{white}]{\textbullet}%
  \,\textsc{\sffamily\small\color{accent}Finding}}, #1}

% --- Key question: structurally distinct -- full thin frame, not left-rule
\newtcolorbox{keyquestionbox}[1][]{
  enhanced, breakable,
  colback=white, coltext=ink,
  colframe=accent, boxrule=0.6pt,
  arc=1pt, left=9pt, right=9pt, top=6pt, bottom=6pt,
  before skip=8pt, after skip=8pt,
  fonttitle=\sffamily\bfseries\small,
  attach title to upper={\par\vspace{2pt}\noindent},
  title={\tcbox[on line, colback=white, colframe=accent, boxrule=0.6pt, arc=1pt,
    left=2pt, right=2pt, top=0.3pt, bottom=0.3pt,
    fontupper=\sffamily\bfseries\fontsize{6.3}{6.3}\selectfont\color{accent}]{?}%
  \,\textsc{\sffamily\small\color{accent}Key question}}, #1}

% --- "How to read this report" front-matter box ----------------------------
\newtcolorbox{howtoreadbox}{
  enhanced, breakable,
  colback=gray100, coltext=ink,
  colframe=gray400, boxrule=0.4pt,
  arc=1pt, left=10pt, right=10pt, top=8pt, bottom=8pt,
  fonttitle=\sffamily\bfseries,
  title={\sffamily\bfseries\color{ink}How to read this report},
}

% --- custom "part" divider: number + rule + title, full control -----------
\newcounter{reportpart}
\newcommand{\reportpart}[2]{%
  \clearpage
  \thispagestyle{empty}
  \phantomsection
  \refstepcounter{reportpart}
  \addcontentsline{toc}{part}{\protect\numberline{\Roman{reportpart}}#1}
  \vspace*{0.12\textheight}
  {\sffamily\bfseries\fontsize{14}{16}\selectfont\color{gray700}PART~\Roman{reportpart}}\par
  \vspace{2mm}
  \noindent\textcolor{accent}{\rule{\linewidth}{1.1pt}}\par
  \vspace{4mm}
  {\sffamily\bfseries\fontsize{28}{32}\selectfont\color{ink}#1}\par
  \ifx\relax#2\relax\else
    \vspace{3mm}
    {\sffamily\fontsize{12}{15}\selectfont\color{gray700}#2}\par
  \fi
  \vspace{0.08\textheight}
}

% ------------------------------------------------ repository references
% \repo{path/in/repo} prints the path in mono and links it to the file on GitHub.
% \code{identifier} prints a code identifier in mono with line-break points.
\usepackage{xurl}
% \auditcommit is the commit whose artifacts the reports were checked against; links are pinned to it.
\newcommand{\auditcommit}{AUDITSHA}
\newcommand{\auditshort}{AUDITSHORT}
\newcommand{\reporelease}{report-v1.1}
\newcommand{\repourl}{https://github.com/AdamKrysztopa/edu_proj/blob/\auditcommit/}
\newcommand{\repo}[1]{\href{\repourl#1}{\nolinkurl{#1}}}
\newcommand{\code}[1]{\texttt{\small\detokenize{#1}}}

% Footnotes hold long repository paths; ragged right avoids stretched spacing.
\deffootnote[1.2em]{1.2em}{1em}{\textsuperscript{\thefootnotemark}\,}
\addtokomafont{footnote}{\raggedright}
